\documentclass{report}
\usepackage{amsmath,amssymb}
\setlength{\parindent}{0mm}
\setlength{\parskip}{1em}
\begin{document}
\begin{center}
\rule{6in}{1pt} \
{\large Gabriel Wittum \\
{\bf UG 4 - a novel multigrid software system for the simulation of pde-based models.}}

G-CSC \\ Universit�t Frankfurt \\ Kettenhofweg 139 \\ 60325 Frankfurt am Main \\ Germany
\\
{\tt wittum@g-csc.de}\end{center}

Numerical simulation has become one of the major topics in Computational
Science. To promote modelling and simulation of complex problems new
strategies are needed allowing for the solution of large, complex model
systems. Crucial issues for such strategies are reliability, efficiency,
robustness, usability, and versatility.

After discussing the needs of large-scale simulation we point out basic
simulation strategies such as adaptivity, parallelism and multigrid
solvers. To allow adaptive, parallel computations the load balancing
problem for dynamically chaniging grids has to be solved efficiently by
fast heuristics. These strategies are combined in the novel simulation
system UG 4 (�Unstructured Grids�) being presented in the following. In
addition to these methodological features, a new GUI concept for reducing
complexity in handling and controlling the simulation and the software is
presented.

In the second part of the talk we show the performance and efficiency of
this strategy in various applications. In particular large scale parallel
computations of density-driven groundwater flow as well as some
non-standard problems from biotechnology and medicine are discussed in
more detail. We present a new model for biomass fermentation and for
signal processing in neurons in three space dimensions.


\end{document}
